\documentclass[a4paper,12pt]{article}
\usepackage[italian]{babel}
\usepackage[T1]{fontenc}
\usepackage[utf8]{inputenc}
\usepackage{geometry}
\geometry{margin=2.5cm}

\title{Esame di basi di dati - Prima parte}
\author{ITS Academy - Fabbrica Digitale 4.0}
\date{Prof. Fedeli Massimo}

\begin{document}
	
	\maketitle
	
	Una clinica privata intende realizzare un sistema informativo per la gestione dei pazienti, dei medici, delle visite mediche e dei reparti presenti nella struttura.
	
	La clinica è organizzata in diversi reparti (ad esempio: cardiologia, ortopedia, neurologia, pediatria, radiologia). Per ogni reparto devono essere memorizzati: codice reparto, nome reparto, piano dell’edificio e numero di posti disponibili.
	
	I medici lavorano presso la clinica e sono assegnati a un reparto. Per ogni medico devono essere memorizzati: codice medico, nome, cognome, specializzazione, telefono, email e reparto di appartenenza. Un reparto può avere più medici, mentre ogni medico appartiene a un solo reparto.
	
	I pazienti possono prenotare visite mediche. Per ogni paziente devono essere memorizzati: codice paziente, nome, cognome, data di nascita, codice fiscale, indirizzo, telefono, email.
	
	Le visite mediche rappresentano gli appuntamenti tra pazienti e medici. Per ogni visita devono essere memorizzati: codice visita, paziente, medico, data visita, ora visita, tipo di visita, diagnosi, eventuale terapia prescritta e costo della visita.
	
	Un paziente può effettuare più visite nel tempo con medici diversi, mentre ogni visita è associata a un solo paziente e a un solo medico. Un medico può effettuare molte visite.
	
	La clinica dispone anche di alcune sale visita. Per ogni sala devono essere memorizzati: codice sala, numero sala, piano e reparto di appartenenza. Ogni visita si svolge in una determinata sala.
	
	Il sistema deve permettere di conoscere:
	\begin{itemize}
		\item lo storico delle visite di ogni paziente;
		\item l’elenco delle visite effettuate da ogni medico;
		\item l’occupazione delle sale visita in una certa data;
		\item il numero di visite effettuate per ogni reparto.
	\end{itemize}
	
	\bigskip
	
	\textbf{Richieste}
	
	Sulla base della descrizione fornita si richiede di:
	
	\begin{enumerate}
		\item Individuare le entità principali del dominio, i loro attributi e le chiavi primarie.
		\item Definire le relazioni tra le entità, specificandone cardinalità e partecipazione (obbligatoria o facoltativa).
		\item Disegnare lo schema E-R.
		\item Tradurre lo schema E-R in schema relazionale.
	\end{enumerate}
	
	Consegna: inviare al docente tramite email lo schema E-R con nome file "cognome nome.json"
	e lo schema logico all'interno di un file pdf denominato "logicocognomenome.pdf".
	
\end{document}